\title{Is Reinforcement Learning (Not) for Natural Language Processing: Benchmarks, Baselines, and Building Blocks for Natural Language Policy Optimization}

\begin{abstract}
{\let\thefootnote\relax\footnotetext{$^\textbf{*}$Denotes Equal Contribution}} We tackle the problem of aligning pre-trained large language models (LMs) with human preferences. If we view text generation as a sequential decision-making problem, reinforcement learning (RL) appears to be a natural conceptual framework. However, using RL for LM-based generation faces empirical challenges, including training instability due to the combinatorial action space, as well as a lack of open-source libraries and benchmarks customized for LM alignment. Thus, a question rises in the research community:
\emph{is RL a practical paradigm for NLP?} 

To help answer this, we first introduce an open-source modular library,
\textbf{RL4LMs\footnote{Code: \url{https://github.com/allenai/RL4LMs}},\footnote{Project Website: \url{https://rl4lms.apps.allenai.org/}}} 
for optimizing language generators with RL. The library consists of on-policy RL algorithms that can be used to train any encoder or encoder-decoder LM in the HuggingFace library \citep{wolf2020transformers} with an arbitrary reward function. Next, we present the \textbf{GRUE (General Reinforced-language Understanding Evaluation)} benchmark, a set of 6 language generation tasks which are supervised not by target strings, but by reward functions which capture automated measures of human preference. GRUE is the first leaderboard-style evaluation of RL algorithms for NLP tasks. Finally, we introduce an easy-to-use, performant RL algorithm, \textbf{\algorithmfont{NLPO} (Natural Language Policy Optimization)} that learns to effectively reduce the combinatorial action space in language generation. We show 1) that RL techniques are generally better than supervised methods at aligning LMs to human preferences; and 2) that NLPO exhibits greater stability and performance than previous policy gradient methods (e.g., PPO~\citep{schulman2017proximal}), based on both automatic and human evaluations.
\end{abstract}

\section{Introduction} \label{sec:main/10_intro} \input{tex/main/10_intro}
\section{Related Work} \label{sec:main/20_background} \input{tex/main/20_background}
\section{RL4LMs: A Library for Training LMs with RL} \label{sec:main/30_framework} \input{tex/main/30_framework}
\section{NLPO: Natural Language Policy Optimization} \label{sec:main/40_nlpo} \input{tex/main/40_nlpo}
\section{GRUE (General Reinforced-language Understanding Eval)} \label{sec:main/50_evaluation} \input{tex/main/50_evaluation}
\section{Conclusions} \label{sec:main/60_conclusion} \input{tex/main/60_conclusion}

\end{document}
